\documentclass[12pt]{article}

\usepackage[utf8]{inputenc}
\usepackage[romanian]{babel}
\usepackage{amsmath, amssymb, amsthm}
\usepackage{geometry}
\geometry{a4paper, margin=2.5cm}

\title{Analiză Funcțională — Spațiu Separabil}
\author{}
\date{}

\theoremstyle{definition}
\newtheorem{definition}{Definiție}

\begin{document}

\maketitle

\begin{definition}
Un spațiu metric $(X,d)$ se numește \emph{separabil} dacă există o mulțime numărabilă
$D \subseteq X$ care este densă în $X$, adică


\[
\overline{D} = X.
\]


\end{definition}

\begin{definition}
Un spațiu normat $(E,\|\cdot\|)$ se numește \emph{separabil} dacă spațiul metric
indus de normă este separabil, adică există o mulțime numărabilă $D \subseteq E$
cu proprietatea că


\[
\forall x \in E,\ \forall \varepsilon > 0,\ \exists d \in D \text{ astfel încât }
\|x-d\| < \varepsilon.
\]


\end{definition}

\end{document}
